% SGA 3, Expose VI_B, Loop 1 English reconstruction.
% Scope: Remark 6.3, Proposition 6.4, Lemmas 6.4.2--6.4.3,
% and Definition 6.5 in full, including source editor's note 60.
% Authority: Polo--Gille Exp6B-13oct24.pdf, component-local pages
% 44--45 (combined-reader pages 385--386; re-edition headers
% 374--375; original-margin printed pages 359--360).
% Proposition 6.5.1 begins on local page 45 / combined page 386
% and is excluded.
% OCR and Reinhold English Markdown are comparison witnesses only.
% This unit contains no diagrams.

\medskip
\phantomsection\label{rem:sga3-VIB-6.3}%
\noindent\textbf{注 6.3.---}
\emph{设 \(k\) 为代数闭域, \(G\) 为
\(k\)-群, \(H\) 为 \(G\) 的子群概形. 假设
\(G\) 和 \(H\) 在 \(k\) 上有限型且
约化. 则
\(\underline{\operatorname{Norm}}_G H\) (相应地,
\(\underline{\operatorname{Centr}}_G H\)) 可由
\(G\) 的子群概形表出, 其伴随约化子概形
恰为 \emph{Bible} 意义下 \(H\) 在
\(G\) 中的正规化子 (相应地, 中心化子).}

\medskip
\phantomsection\label{prop:sga3-VIB-6.4}%
\noindent\textbf{命题 6.4.---}
\emph{设 \(G\) 为 \(S\)-群, \(u:X\to G\) 为
\(S\)-概形的单态射. 令}
\[
  T=\underline{\operatorname{Transp}}_G(X,X).
\]
\emph{下列条件等价:}
\begin{enumerate}
\item[\textup{(i)}]\phantomsection\label{cond:sga3-VIB-6.4-i}
  \emph{\(T\) 是 \(G\) 的群子 \(S\)-函子.}
\item[\textup{(ii)}]\phantomsection\label{cond:sga3-VIB-6.4-ii}
  \emph{\[
    T=\underline{\operatorname{Transp}}^{\mathrm{str}}_G(X,X)
     =\underline{\operatorname{Norm}}_G X.
  \]}
\end{enumerate}
\emph{在下列两种情形中的每一种下, 这些条件均成立:}
\begin{enumerate}
\item[\textup{a)}]\phantomsection\label{case:sga3-VIB-6.4-a}
  \emph{\(X\) 在 \(S\) 上有限呈示.}
\item[\textup{b)}]\phantomsection\label{case:sga3-VIB-6.4-b}
  \emph{\(T\) 可由在 \(S\) 上有限
  呈示的概形表出.}
\end{enumerate}

\textup{(i)} 和 \textup{(ii)} 的等价性由如下事实得到:
对每个态射 \(S'\to S\) 及任意
\(t,t'\in T(S')\), 都有
\[
  tt'\in T(S'),
\]
并且有恒等式
\[
  \underline{\operatorname{Transp}}^{\mathrm{str}}_G(X,X)
  =T\cap c(T)
\]
(参见 \hyperref[def:sga3-VIB-6.1]{6.1\textup{(ii)}}).

先假设 \textup{a)} 成立. 若 \(t\in T(S)\), 则
\(\operatorname{int}(t)\) 为 \(X\) 到自身的单态射,
因而由 EGA~IV\(_4\), 17.9.6, 它是 \(X\) 的 \(S\)-自同构.
因此
\[
  t\in\underline{\operatorname{Transp}}^{\mathrm{str}}_G(X,X),
\]
这就证明了情形 \textup{a)} 下的断言.

在情形 \textup{b)} 中, 显然有
\[
  \mu(T\times_S T)\subset T,
\]
断言由下面的引理得到.

\medskip
\phantomsection\label{lem:sga3-VIB-6.4.2}%
\noindent\textbf{引理 6.4.2.---}%
\footnote{\label{note:sga3-VIB-60}编者注: 我们保留了
原文的编号; 没有编号 6.4.1.}
\emph{设 \(G\) 为有限呈示的 \(S\)-概形,
赋有一个结合运算 (EGA~0\(_{\mathrm{III}}\), 8.2.5
意义下). 假设对每个
\(S\)-概形 \(S'\) 及每个 \(g\in G(S')\), 在集合 \(G(S')\) 中
由 \(g\) 定义的右平移和左平移均为单射,
并假设 \(G(S)\ne\varnothing\). 则 \(G\) 为
\(S\)-群.}

\emph{证明.}
只需证明对每个 \(S\)-概形 \(S'\), \(G(S')\) 为群.
假设立即蕴涵, 对每个 \(g\in G(S')\),
在 \(G_{S'}\) 上由其定义的右平移和左平移
均为 \(S'\)-单态射
\[
  G_{S'}\longrightarrow G_{S'}.
\]
因此它们都是 \(S'\)-自同构, 因为 \(G\) 在
\(S\) 上有限呈示 (EGA~IV\(_4\), 17.9.6).
因而在集合 \(G(S')\) 中由 \(g\) 定义的
右平移和左平移均为双射, 断言化归为
下面的引理.

\medskip
\phantomsection\label{lem:sga3-VIB-6.4.3}%
\noindent\textbf{引理 6.4.3.---}
\emph{设 \(G\) 为带有结合运算的非空集合,
使得右平移和左平移均为双射. 则 \(G\) 为
群.}

证明留给读者.

\medskip
\phantomsection\label{def:sga3-VIB-6.5}%
\noindent\textbf{定义 6.5.---}
\emph{设 \(G\) 为 \(S\)-群, \(H\) 为
\(S\)-函子, \(u:H\to G\) 为单态射.}
\begin{enumerate}
\item[\textup{(i)}]
  \emph{\(H\) 在 \(G\) 中的连通中心化子}, 记为
  \[
    \underline{\operatorname{Centr}}^0_G H,
  \]
  是函子
  \(\underline{\operatorname{Centr}}_G H\) 的单位连通分支
  (参见 \hyperref[def:sga3-VIB-3.1]{3.1} 及
  \hyperref[def:sga3-VIB-6.1]{6.1\textup{(iii)}}).
\item[\textup{(i\(^{\prime}\))}]
  对每个态射 \(u:X\to G\), 类似地定义
  函子
  \[
    \underline{\operatorname{Centr}}^0(u)
  \]
  (参见 \SGARefLink{prop:sga3-VIB-6.2}{6.2}\textup{(iv)}).
\item[\textup{(ii)}]
  当对每个 \(s\in S\), \(u_s\) 都为闭浸入时,
  \emph{\(H\) 在 \(G\) 中的连通正规化子}, 记为
  \[
    \underline{\operatorname{Norm}}^0_G H,
  \]
  是函子
  \(\underline{\operatorname{Norm}}_G H\) 的单位连通分支.
\end{enumerate}

\noindent\emph{注意.}
由 \hyperref[cor:sga3-VIB-1.4.2]{1.4.2}, 当 \(H\) 为 \(S\)-群概形,
\(G\) 和 \(H\) 在 \(S\) 上局部有限型, 且
\(u\) 为 \(S\)-群的拟紧态射时,
\textup{\SGARefLink{cor:sga3-VIB-1.4.2:item:ii}{(ii)}} 中的假设满足.

% 翻译统计: 源行数 152; 输出行数 154 (含本统计行); 行内/行间数学片段 89; 翻译占位 0. 底本疑点另报主线程.
